\section{The Way of Gnosis}

Gnosis is one of the Gifts of the Holy Spirit. It is not secret knowledge available for the few, nor is it the ``real"
interpretation of text, opposed to the orthodox. On the other hand, it is perhaps an Open Secret, since it is pursued
by the few. It is not necessary for salvation, just one path for those so inclined, just as others my respond better to
different gifts.

We are asked to know, love, and serve God. We cannot love what, or whom, we do not know. We cannot server what, or whom,
we do not love. So, in a sense, knowing, or gnosis, is fundamental to the spiritual life.

Gnosis is a method of depth. Hence, it is not merely intellectual, since it also affects the Feeling and Willing aspects
of life.

\paragraph{Gnosis and the Text}
Gnosis is a direct knowing, as opposed to a rationalist or discursive knowing. This is the way it manifests:

Mysticism $\Rightarrow$ Gnosis $\Rightarrow$ Text

It necessarily begins with mystical experience, which is not earned, but is rather a free gift. That leads to gnosis,
which then tries to explain itself through language. The resulting texts get debated while its source – gnosis – is
forgotten and pushed to the background. The process cannot work in reverse. Viz.,

Text $\Rightarrow$ Gnosis $\Rightarrow$ Mysticism

Although texts certainly can be helpful, especially as road markers, all the study and debate over texts cannot induce
gnosis, but often misunderstanding.

\paragraph{Preparation}
It is first necessary to reduce the perturbations in one's soul life in order to reflect the Spirit
more faithfully. These are the main preliminary steps in that process:

\begin{itemize}
\item Concentration without effort 
\item Moral purification of the Will 
\item Overcoming of the subjective element in the Intellect 
\item Non-expression of negative emotions 
\end{itemize}
We can provide some examples.

\paragraph{Christian Initiation}
The three sacraments of Christian initiation can be expressed like this.

\begin{itemize}
\item \textbf{Baptism}: The natural born man is alienated from God. Through Baptism, he is restored and becomes a child
of God by adoption. The popular view that those born of the flesh are children of God is incorrect. Being a child of
God means to share in his nature; hence our task is the achievement of theosis. 
\item \textbf{Communion}: Ideally, communion is administered to children who reach the age of reason around seven years
old. That does not mean that they can suddenly do logic puzzles, but rather that they can distinguish between good and
evil. Thus, the child reverses the choice made by Adam and Eve who wanted to be their own arbiters of right and wrong. 
\item \textbf{Confirmation}: At confirmation we receive the gifts of the Holy Spirit. Most remain passive to those
gifts, accepting them when given. On the path of gnosis, however, one must play an active role in order to make oneself
available for the gifts of knowledge and understanding. 
\end{itemize}
\paragraph{Inner Soul Life}
These are the elements in the philosophical understanding of the soul as described by Thomas Aquinas:

\begin{itemize}
\item \textbf{Sensitive soul}: In common with all living things 
\item \textbf{Animal soul}: In common with animals 
\item \textbf{Intellectual soul}: specific to man and angels 
\item These form a unity, i.e., the person or the I 
\end{itemize}
That is correct philosophically. However, for gnosis, one learns to observe the activities specific to the different
souls and how they relate to each other. One sees that the inner unity is only virtual and not fully actualized. In
other words, there is a war in the soul as long as we are not pure of heart. Our task, then, is develop the I to be
able to master the disparate elements of soul life.

\paragraph{Knowledge of God}
The Thomist philosopher \textbf{Edward Feser} describes Thomas' understanding of God as classical
theism\footnote{\url{https://edwardfeser.blogspot.com/2012/07/classical-theism-roundup.html}}. This was the understanding of the ancients, with these examples provided by Feser:

\begin{itemize}
\item \textbf{Maimonides}: Jewish 
\item \textbf{Avicenna}: Muslim 
\item \textbf{Plotinus}: Pagan 
\end{itemize}
Although we share a common understanding of God with these philosophers, this does not imply a religious syncretism, nor
agreement with different creeds and forms of worship. Yet we cannot deny the same basic understanding of God among
them. Since Feser is a philosopher, he often argues against atheists. Moreover, he also opposes a different God
conception called neotheism or theistic personalism. Not only is this view held by some contemporary theologians, it is
probably the implicit understanding of common believers to one extent or another.

These are the main features held by neotheists:

\begin{itemize}
\item God is a being among other beings 
\item He is subject to time 
\item God has parts and is not simple 
\item He experiences human emotions, that is, He is not impassable 
\item He has a gnomic will and deliberates before making decisions 
\item His knowledge is propositional 
\item He has sensory experiences 
\end{itemize}
However, when one has a direct experience of God, then classical theism is simply understood beyond argument. Atheism
and neotheism are not just wrong, they are incoherent. Now I don't want to call neotheists
heretics nor do I deny that they may provide some important insights into the nature of God. Yet their arguments never
settle any questions.

God's knowledge is not propositional, since His knowledge is direct knowing of essences in the
Divine Mind. When that point is understood, the understanding of Wisdom, or Sophia, arises. Hence the next steps are to
follow the paths described by the Christian Cabala, Jacob Boehme, and Vladimir Solovyov.



\flrightit{Posted on 2019-10-13 by Cologero }

\begin{center}* * *\end{center}

\begin{footnotesize}\begin{sffamily}



\texttt{Fallen Adam on 2019-10-13 at 14:47 said: }

Would our knowledge of the beings in the sensory world count as a low degree of mysticism? Referring here to our act of
knowing what a being is trough the perceptions we obtain of it trough the senses. It's an
intuitive and direct form of knowing, in that any reason we might give to justify such knowledge, is only posterior to
the act itself. We can justify it by saying that we perceived trough the senses is similar to other such group of
perceptions, but we all know that we say ``I saw a tree", we are referring to a being which shares in a certain nature
and essence, and not to a bundle of sense perceptions.

Also, could the experience of observing and understanding what is happening in our inner life be a form of mysticism?
For example, the moment in which we apprehend the fact that we are being moved by anger in a particular situation, or
that negative emotions are triggered in us due to a certain cause. This might also be called a discursive process, in
that we might use our memory and give our self reasons towards reaching a particular conclusion, but there is also an
aspect of immediacy and direct intuition in such things.


\hfill

\texttt{Cassiodorus on 2019-10-13 at 21:12 said: }

I don't see how the classical theism articulated by Feser is reconcilable with the nondualism
espoused by Guenon .


\hfill

\texttt{Cologero on 2019-10-13 at 21:34 said: }

We have already mentioned on various occasions that in principle Aristotle

posited identification by knowledge, but also that this affirmation, in his works as

in those of his Scholastic followers, seems to have remained purely theoretical, for

they seem never to have drawn any conclusions from it as concerns metaphysical

realization (From the Multiple States of the Being).

Classical theism is fine as far as it concerns Being. Guenon refers to non-being. Wolfgang Smith recognizes this
incompleteness and uses Jacob Boehme and Meister Eckhart as the basis for a ``Trinitarian non-duality". Ultimately,
metaphysical realization — what we are calling mysticism and gnosis — is the
final arbiter.


\hfill

\texttt{Sviatoslav on 2019-10-14 at 09:27 said: }

It seems that everything starts with mystical experience. In this context, I guess mysticism has a broader meaning than
Guenon or Evola understood that. Maybe more in Tomberg's version. But is mysticism only way, or one of
the ways? Tomberg distinguishes three forms of mystical experience :

– the experience of union with Nature

– union with the transcendental human Self

– union with God

The first leads to intoxication, the second to sobriety and the third is a synthesis of both previous. 

Are all of those ways valid, or one is better than others? Tomberg exalts union with God in love, in God-man duality
while discourages from intellectual Vedantic path (including Guénon himself). Maybe I am wrong, but it seems to me,
that there is a significant conflict between Guénon and Tomberg.


\hfill

\texttt{Casssiodorus on 2019-10-14 at 11:36 said: }

I just do not see how the doctrine of maya in divinis in which Being is relegated to a lower level of Divinity can be
squared with the classical theistic tradition of Aristotle, Aquinas, or even that of the Dvaitan Vedantan, Madhva
.

Classical theism speaks of the ontological distinction and contingency of creation, insisting that differences are real.


The unqualified nondualism of Guenon, following Advaita Vedanta, rejects the reality of difference and claims that all
that is relative is maya or ``illusion" – including Being Itself which Schuon called the ``relative Absolute".

These principles are fundamentally pantheistic , or panentheistic if you prefer, and run contrary to the doctrine of
creation ex nihilo. Furthermore, the Trinity which is essentially a matter of relationship would also be relegated to
the domain of maya or illusion. Problematic to say the least for the traditional Christian.


\hfill

\texttt{Casssiodorus on 2019-10-14 at 14:51 said: }

Does the distinction between classical theism and theistic personalism relate

to the distinction between ``non-Being" and Being?


\hfill


\end{sffamily}\end{footnotesize}
